% Part:sets-functions-relations
% Chapter: size-of-sets
% Section: pairing

\documentclass[../../../include/open-logic-section]{subfiles}

\begin{document}

\olfileid{sfr}{siz}{pai}
\olsection{যুগলায়ন অপেক্ষক ও সংকেত}

\begin{explain}
কান্টরের আঁকাবাঁকা পথের পদ্ধতি $\Nat^n$-এর তালিকায়নযোগ্যতাকে ছবির মাধ্যমে স্পষ্ট করে। তবে এবার $\Nat^2$-এর সারণিটির দিকে নজর দিই। সারণিতে আঁকাবাঁকা পথ ধরে ক্রমস্থান গুনে যাচাই করতে পারি যে $\tuple{1,2}$-এর সঙ্গে~$7$ সংখ্যাটি যুক্ত। কিন্তু আরও সরাসরি এটি গণনা করতে পারলে ভালো হতো। অর্থাৎ আঁকাবাঁকা তালিকায়নের \emph{বিপরীত} $g\colon \Nat^2 \to \Nat$ আমাদের হাতে থাকলে সুবিধা হতো, যাতে
\[
g(\tuple{0,0}) = 0, \;
g(\tuple{0,1}) = 1, \;
g(\tuple{1,0}) = 2, \; \dots,
g(\tuple{1,2}) = 7, \; \dots
\]
এর সাহায্যে তালিকায়নে $\tuple{n, m}$ ঠিক কোথায় আসবে, তা গণনা করতে পারতাম।

বস্তুত, দুটি বিষয় লক্ষ করে সরাসরি $g$ সংজ্ঞায়িত করতে পারি। প্রথমত, $n$-তম সারি ও $m$-তম স্তম্ভে~$v$ মানটি থাকলে, $(n+1)$-তম সারি ও $(m-1)$-তম স্তম্ভে $v + 1$ থাকে। দ্বিতীয়ত, আমাদের তালিকায়নের প্রথম সারিতে ত্রিভুজসংখ্যাগুলি আছে, যেগুলি $0$, $1$, $3$, $6$ ইত্যাদি দিয়ে শুরু হয়। $k$-তম ত্রিভুজসংখ্যা হলো $< k$ স্বাভাবিক সংখ্যাগুলির যোগফল, যা $k(k+1)/2$ দিয়ে গণনা করা যায়। এই দুটি পর্যবেক্ষণ একত্র করে নিচের অপেক্ষকটি বিবেচনা করো:
\[
  g(n,m) = \frac{(n+m+1)(n+m)}{2} + n
\]
পড়তে সুবিধা হয় বলে আমরা প্রায়ই $g(\tuple{n, m})$-এর পরিবর্তে শুধু $g(n, m)$ লিখি। এই সূত্র অনুযায়ী প্রথমে $(n+m)^\text{তম}$ ত্রিভুজসংখ্যা নির্ণয় করতে হবে, তারপর তার সঙ্গে $n$ যোগ করতে হবে। এতে সারণিটি ঠিক আমাদের কাঙ্ক্ষিতভাবে পূর্ণ হয়। বিশেষত ক্রমযুগল $\tuple{1, 2}$, $\frac{4 \times 3}{2} +
1 = 7$-এ পাঠানো হয়।

এই অপেক্ষক $g$ একটি ক্রমযুগলের সেটের তালিকায়নের \emph{বিপরীত}। এই ধরনের অপেক্ষককে \emph{যুগলায়ন অপেক্ষক} বলা হয়।
\end{explain}

\begin{defn}[যুগলায়ন অপেক্ষক]
  একটি অপেক্ষক $f\colon A \times B \to \Nat$-এর ক্ষেত্রে $f$ একৈক হলে তাকে পাটিগাণিতিক \emph{যুগলায়ন অপেক্ষক} বলা হয়। আমরা আরও বলি যে $f$, $A \times B$-এর \emph{সংকেতায়ন} করে এবং $f(x,y)$ হলো $\tuple{x,y}$-এর \emph{সংকেত}।
\end{defn}

\begin{explain}
যুগলায়ন অপেক্ষক দিয়ে, যেমন, স্বাভাবিক সংখ্যার ক্রমযুগলের সংকেতায়ন করতে পারি; অন্যভাবে বললে, উপাদানের প্রত্যেক \emph{যুগলকে} \emph{একটিমাত্র} সংখ্যা দিয়ে প্রকাশ করতে পারি। যুগলায়ন অপেক্ষকের বিপরীত ব্যবহার করে সংখ্যাটির \emph{সংকেতোদ্ধার} করতে পারি, অর্থাৎ সেটি কোন যুগলকে প্রকাশ করে তা বার করতে পারি।
\end{explain}

\begin{prob}
সকল অঋণাত্মক মূলদ সংখ্যার সেটের একটি তালিকায়ন দাও।
\end{prob}

\begin{prob}
দেখাও যে $\Rat$ !!{enumerable}। মনে রেখো, যেকোনো মূলদ সংখ্যাকে $z/m$ ভগ্নাংশ হিসেবে লেখা যায়, যেখানে $z \in \Int$, $m \in \Nat^+$।
\end{prob}

\begin{prob}
$\Bin^*$-এর একটি তালিকায়ন সংজ্ঞায়িত করো।
\end{prob}

\begin{prob}
যুক্তিবিদ্যার প্রাথমিক পাঠ থেকে মনে করো, প্রত্যেক সম্ভাব্য সত্যসারণি একটি সত্যমান-অপেক্ষক প্রকাশ করে। অন্যভাবে বললে, সত্যমান-অপেক্ষকগুলি হলো কোনো~$k$-এর জন্য $\Bin^k \to \Bin$-এর সকল অপেক্ষক। প্রমাণ করো যে সকল সত্যমান-অপেক্ষকের সেট তালিকায়নযোগ্য।
\end{prob}

\begin{prob}
দেখাও যে যেকোনো অসীম !!{enumerable} সেটের সকল সসীম উপসেটের সেট !!{enumerable}।
\end{prob}

\begin{prob}
$\Nat$-এর কোনো উপসেটকে \emph{সহসসীম} বলা হয় যদি এবং কেবল যদি সেটি $\Nat$-এর মধ্যে কোনো সসীম সেটের পূরক হয়; অর্থাৎ $A \subseteq \Nat$ সহসসীম হয় যদি এবং কেবল যদি $\Nat\setminus A$ সসীম হয়। ধরি, $I$ সেই সেট যার !!{element}s ঠিক $\Nat$-এর সসীম ও সহসসীম উপসেটগুলি। দেখাও যে $I$ !!{enumerable}। % OLSIZ-002 TARGET-CORRECTION-BEGIN
\par\smallskip\noindent\textnormal{\small\textbf{উৎস-ব্যাখ্যা (OLSIZ-002):} হিমায়িত ইংরেজি উৎসের ``complement of a finite set $\Nat$'' বাক্যাংশটি ব্যাকরণগতভাবে অসম্পূর্ণ। পরের সমতুল্যতা অর্থটি নির্ধারণ করে: $\Nat$-এর কোনো সসীম উপসেটের $\Nat$-এর মধ্যে পূরক। বাংলা পাঠে সেই নির্ধারিত অর্থটি স্পষ্ট করে বলা হয়েছে।} % OLSIZ-002 TARGET-CORRECTION-END
\end{prob}

\begin{prob}
দেখাও যে !!{enumerable} পরিবারে থাকা !!{enumerable} সেটগুলির সংযোগ !!{enumerable}। অর্থাৎ $A_1$, $A_2$, \dots{} সেট হলে এবং প্রত্যেক $A_i$ !!{enumerable} হলে, তাদের সকলের সংযোগ $\bigcup_{i=1}^\infty
A_i$-ও !!{enumerable}। [দ্রষ্টব্য: এটি কঠিন!]
\end{prob}

\begin{prob}
ধরি, $f \colon A \times B \to \Nat$ যেকোনো একটি যুগলায়ন অপেক্ষক। দেখাও যে $f$-এর বিপরীত $A \times B$-এর একটি তালিকায়ন।
\end{prob}

\begin{prob}
$\Nat^3$-এর সংকেতায়ন করে এমন একটি অপেক্ষক নির্দিষ্ট করো।
\end{prob}

\end{document}
